% =====================================================================
%  Annexe B — Références et webographie
% =====================================================================
\chapter{Références et webographie}
\label{ann:references}

Les références ci-dessous correspondent aux documentations effectivement consultées
pendant la réalisation, ainsi qu'aux cadres méthodologiques mentionnés dans le rapport.
Les dates de consultation sont à compléter si l'établissement l'exige.

\section*{Documentations techniques}

\begin{enumerate}[label={[\arabic*]}, leftmargin=*, itemsep=5pt]
  \item \textbf{FastAPI} — documentation officielle du framework, notamment les
    chapitres sur l'injection de dépendances et la sécurité.\\
    \url{https://fastapi.tiangolo.com/}

  \item \textbf{SQLAlchemy 2.0} — guide de l'ORM en style déclaratif typé.\\
    \url{https://docs.sqlalchemy.org/en/20/}

  \item \textbf{Alembic} — documentation des migrations de schéma.\\
    \url{https://alembic.sqlalchemy.org/}

  \item \textbf{Pydantic v2} — validation et sérialisation des données.\\
    \url{https://docs.pydantic.dev/latest/}

  \item \textbf{PostgreSQL} — manuel de référence.\\
    \url{https://www.postgresql.org/docs/}

  \item \textbf{React} — documentation officielle.\\
    \url{https://react.dev/}

  \item \textbf{TypeScript} — manuel du langage.\\
    \url{https://www.typescriptlang.org/docs/}

  \item \textbf{Vite} — documentation de l'outil de build.\\
    \url{https://vite.dev/}

  \item \textbf{React Router} — routage côté client.\\
    \url{https://reactrouter.com/}

  \item \textbf{TanStack Query} — gestion du cache de données serveur.\\
    \url{https://tanstack.com/query/latest}

  \item \textbf{Recharts} — bibliothèque de graphiques pour React.\\
    \url{https://recharts.org/}

  \item \textbf{pytest} — documentation du cadre de test.\\
    \url{https://docs.pytest.org/}

  \item \textbf{Docker Compose} — référence du fichier de composition.\\
    \url{https://docs.docker.com/compose/}
\end{enumerate}

\section*{Normes et spécifications}

\begin{enumerate}[label={[\arabic*]}, leftmargin=*, itemsep=5pt, start=14]
  \item \textbf{RFC 7519} — \emph{JSON Web Token (JWT)}, Internet Engineering Task
    Force, mai 2015. Spécification du format de jeton utilisé pour
    l'authentification.\\
    \url{https://datatracker.ietf.org/doc/html/rfc7519}

  \item \textbf{RFC 6749} — \emph{The OAuth 2.0 Authorization Framework}, IETF, octobre
    2012. Cadre dont FastAPI reprend le flux par mot de passe pour le formulaire de
    connexion.\\
    \url{https://datatracker.ietf.org/doc/html/rfc6749}

  \item \textbf{Spécification OpenAPI} — format de description d'API généré
    automatiquement par l'application sur \code{/docs}.\\
    \url{https://spec.openapis.org/}

  \item \textbf{bcrypt} — Provos, N. et Mazières, D., \emph{A Future-Adaptable Password
    Scheme}, USENIX Annual Technical Conference, 1999. Algorithme de hachage employé
    pour les mots de passe.
\end{enumerate}

\section*{Cadres méthodologiques et référentiels métier}

\begin{enumerate}[label={[\arabic*]}, leftmargin=*, itemsep=5pt, start=18]
  \item \textbf{TOGAF} — \emph{The Open Group Architecture Framework}, The Open Group.
    Cadre dont s'inspirent le découpage du diagnostic AS-IS et le séquencement des
    phases de projet employés dans la plateforme.\\
    \url{https://www.opengroup.org/togaf}

  \item \textbf{Bonnes pratiques GxP} — ensemble des bonnes pratiques applicables à
    l'industrie pharmaceutique (fabrication, laboratoire, distribution), qui fondent
    les garde-fous ajoutés automatiquement par le moteur de qualification lorsqu'une
    opportunité est déclarée en périmètre réglementé.

  \item \textbf{Principes ALCOA+} — exigences d'intégrité des données réglementées
    (attribuable, lisible, contemporain, original, exact, complet, cohérent, durable,
    disponible). Ces principes justifient le caractère immuable du journal d'audit.
\end{enumerate}

\begin{remarque}[Sur la portée de cette liste]
Cette webographie recense les sources techniques réellement utilisées. Les références
méthodologiques TOGAF, GxP et ALCOA+ ont servi de cadre conceptuel au découpage
fonctionnel ; la plateforme s'en inspire sans prétendre à une conformité certifiée,
laquelle relèverait d'une démarche de validation formelle non menée dans ce projet.
\end{remarque}
